\chapter{Vascular Surgery \& Angiology}

\medterm{Acute Limb Ischemia (ALI)}-- Sudden, rapid decrease in lower extremity arterial perfusion threatening limb viability, typically caused by acute arterial embolism (AFib) or in situ thrombosis of an atherosclerotic vessel. Clinical hallmark is the '6 Ps': Pain, Pallor, Pulselessness, Paresthesia, Paralysis, and Poikilothermia. Requires emergency revascularization (surgical embolectomy, catheter-directed thrombolysis, or bypass) within 6 hours to prevent irreversible ischemic neuromuscular damage and amputation.

\medterm{Ankle-Brachial Index (ABI)}-- Non-invasive bedside diagnostic test for peripheral artery disease (PAD), calculated by dividing the highest systolic blood pressure at the ankle (dorsalis pedis or posterior tibial artery) by the highest brachial systolic blood pressure. Interpretation: Normal ($1.00\text{--}1.40$), Borderline ($0.91\text{--}0.99$), Abnormal/PAD ($\le 0.90$), Severe PAD ($<0.50$, associated with critical limb ischemia), and Non-compressible/Calcinosis ($>1.40$, seen in severe diabetes or end-stage renal disease due to arterial medial calcification).

\medterm{Aortic Dissection}-- Life-threatening vascular emergency occurring when a tear in the aortic intima allows blood to surge into the aortic media, creating a false lumen that propagates along the aorta. Classified by Stanford system: Type A (involves ascending aorta, requires emergency open surgical repair) and Type B (confined to descending aorta distal to left subclavian origin, managed medically with blood pressure control or endovascular stent repair—TEVAR). Risk factors: severe chronic hypertension, Marfan syndrome, Ehlers-Danlos syndrome, bicuspid aortic valve, and cocaine use. Presentation: sudden onset of severe, knife-like, tearing or ripping chest pain radiating to the interscapular back. Complications: pericardial tamponade, aortic regurgitation, stroke, and visceral/limb malperfusion.

\medterm{Arteriovenous Fistula (AVF)}-- Abnormal direct communication between an artery and a vein, bypassing the high-resistance capillary bed. Congenital AVFs result from developmental vascular malformations. Acquired AVFs are created surgically (e.g., radiocephalic or brachiocephalic AV fistula in the forearm/arm) as a durable vascular access site for maintenance hemodialysis in end-stage renal disease, or develop traumatically from penetrating vascular injury. Hemodynamic consequences: low resistance causes high flow, palpable thrill, continuous machinery murmur, elevated venous pressure, and potential high-output heart failure or steal syndrome (ischemia of distal tissue).

\medterm{Arteriovenous Graft (AVG)}-- Surgical vascular access for hemodialysis created by interposition of a synthetic conduit (PTFE graft) connecting a peripheral artery (e.g., brachial artery) to a adjacent vein (e.g., basilic vein) when native superficial veins are inadequate for a direct AV fistula.

\medterm{Atherectomy}-- Endovascular interventional technique utilized to mechanically cut, shave, sand (rotational), or vaporize (laser) calcified or heavy atherosclerotic plaque from the lumen of peripheral arteries (e.g., superficial femoral artery) prior to balloon angioplasty or stenting.

\medterm{Axillofemoral Bypass}-- Extra-anatomic arterial bypass procedure performing a subcutaneous graft connection from the axillary artery to the femoral artery; indicated in high-risk patients with severe aortoiliac occlusive disease or infected aortic prosthetic grafts who cannot tolerate major transabdominal aortic reconstruction.

\medterm{Buerger Disease (Thromboangiitis Obliterans)}-- Non-atherosclerotic, segmental, inflammatory vasculitis affecting small- and medium-sized arteries and veins of the extremities, occurring almost exclusively in young male heavy smokers ($<45\text{ years}$). Leads to digital ischemia, severe pain, distal claudication, and painful ulcerations/gangrene. Absolute smoking cessation is the only effective intervention.

\medterm{Carotid Artery Stenting (CAS)}-- Endovascular alternative to open carotid endarterectomy, involving transfemoral or transcarotid deployment of a self-expanding metallic stent across a carotid artery stenosis under distal embolic protection device guidance; preferred in patients with high surgical risk, prior neck radiation, or recurrent post-endarterectomy restenosis.

\medterm{Carotid Endarterectomy (CEA)}-- Surgical procedure performed to remove atherosclerotic plaque from the carotid artery bifurcation and internal carotid artery, reducing the risk of ischemic stroke. Indications: symptomatic carotid stenosis ($>50\text{--}70\%$ luminal narrowing following TIA or non-disabling stroke) and select asymptomatic stenosis ($>70\text{--}80\%$). Operative technique: longitudinal arteriotomy, meticulous plaque excision, and patch closure (dacron or vein graft) to prevent restenosis. Complications: perioperative stroke, cranial nerve injury (hypoglossal, vagus/recurrent laryngeal, marginal mandibular), hyperperfusion syndrome, and wound hematoma.

\medterm{Carotid Stenosis}-- Atherosclerotic narrowing of the carotid artery, primarily localized at the carotid bifurcation and proximal internal carotid artery. Responsible for ~15--20\% of all ischemic strokes due to artery-to-artery thromboembolism or hypoperfusion. Symptomatic presentation: transient ischemic attack (TIA), amaurosis fugax (transient monocular blindness described as a "curtain coming down over one eye"), or acute ischemic stroke. Diagnostic modalities: Carotid Duplex Ultrasound (initial screening), CT Angiography (CTA), or MR Angiography (MRA). Management: intensive medical therapy (antiplatelets, high-intensity statin, blood pressure control) combined with revascularization (Carotid Endarterectomy or Carotid Artery Stenting—CAS) for significant stenosis.

\medterm{Catheter-Directed Thrombolysis (CDT)}-- Minimally invasive endovascular therapy delivering localized thrombolytic agents (tPA, alteplase) directly into a clot via a multi-sidehole infusion catheter placed fluoroscopically within a thrombosed artery or deep vein (e.g., iliofemoral DVT or acute limb ischemia).

\medterm{Central Venous Stenosis}-- Pathological narrowing of central thoracic veins (subclavian, brachiocephalic, or superior vena cava) commonly resulting from previous central venous catheterization, pacemaker leads, or high-flow hemodialysis access. Presents with ipsilateral upper extremity edema, venous engorgement, and hemodialysis access dysfunction.

\medterm{Chronic Venous Insufficiency (CVI)}-- Pathological condition resulting from incompetent venous valves, deep vein thrombosis, or venous outflow obstruction, leading to impaired venous return and sustained ambulatory venous hypertension in the lower extremities. Clinical manifestations (CEAP classification C1--C6): telangiectasias, varicose veins, lower extremity edema, stasis dermatitis (hyperpigmentation from hemosiderin deposition), lipodermatosclerosis (induration of subcutaneous tissue causing "inverted champagne bottle" leg), and non-healing venous stasis ulcers (typically located superior to the medial malleolus). Management: compression therapy (graduated compression stockings), leg elevation, exercise, endovenous thermal ablation (radiofrequency or laser) for superficial reflux, and wound care.

\medterm{Compartment Syndrome (Vascular)}-- Surgical emergency occurring when elevated intracompartmental tissue pressure within a closed osteofascial compartment exceeds capillary perfusion pressure, leading to tissue ischemia and necrosis. Frequently occurs following reperfusion of an acutely ischemic limb. Requires urgent surgical emergency decompressive Fasciotomy.

\medterm{Critical Limb Ischemia (CLI)}-- Advanced, severe stage of peripheral artery disease (PAD) characterized by chronic ischemic rest pain, non-healing ischemic ulcers, or frank gangrene attributable to objectively proven arterial occlusive disease. Rest pain typically affects the distal foot/toes, worsens at night or when elevating the leg, and is transiently relieved by hanging the foot over the edge of the bed. ABI is usually $<0.40$, with ankle systolic pressure $<50\text{ mmHg}$. Represents an imminent threat of limb loss; requires urgent diagnostic angiography and revascularization (endovascular angioplasty/stenting or surgical bypass) to prevent amputation.

\medterm{Deep Vein Thrombosis (DVT)}-- Formation of a blood clot (thrombus) within a deep vein, most commonly in the deep veins of the lower extremities (popliteal, femoral, iliac veins). Pathophysiology described by Virchow's Triad: (1) Venous stasis (immobility, surgery, travel), (2) Endothelial injury (trauma, catheters), and (3) Hypercoagulability (cancer, factor V Leiden, pregnancy, OCPs). Clinical features: unilateral leg pain, swelling, warmth, erythema, and calf tenderness. Complications: Pulmonary Embolism (PE—life-threatening detachment of clot traveling to pulmonary arteries) and Post-Thrombotic Syndrome (PTS). Diagnosis: D-dimer (high sensitivity, excludes DVT), Compression Ultrasonography (gold standard). Treatment: anticoagulation (DOACs, LMWH, warfarin) for at least 3--6 months.

\medterm{Endovascular Aneurysm Repair (EVAR)}-- Minimally invasive surgical technique for repairing abdominal aortic aneurysms. Under fluoroscopic guidance, a modular bifurcated synthetic stent-graft is delivered via transfemoral catheter access and deployed inside the aneurysmal sac, relining the aorta and excluding the aneurysm from systemic arterial pressure. Advantages: lower perioperative morbidity, shorter hospital stay, and lower 30-day mortality compared to open repair. Limitations: requires favorable aortic anatomy (adequate proximal landing zone neck) and necessitates lifelong imaging surveillance (CT or ultrasound) to monitor for endoleaks (persistent blood flow into the aneurysm sac outside the stent-graft).

\medterm{Endovenous Laser Therapy (EVLT)}-- Minimally invasive thermal ablation procedure for superficial venous reflux (saphenous insufficiency); a laser fiber inserted under ultrasound guidance delivers thermal energy to wall of vein, causing collapse, endothelial destruction, and fibrotic occlusion of the incompetent saphenous vein.

\medterm{Fasciotomy}-- Emergency surgical procedure consisting of longitudinal incisions through skin and underlying deep fascia to relieve abnormally high intracompartmental pressure, used to treat acute compartment syndrome following trauma or revascularization of acute limb ischemia.

\medterm{Femorofemoral Bypass}-- Extra-anatomic arterial bypass graft constructed subcutaneously across the suprapubic space, diverting blood from a healthy donor femoral artery to the contralateral ischemic femoral artery in patients with unilateral iliac artery occlusion.

\medterm{Femoropopliteal Bypass}-- Surgical arterial revascularization procedure bypassing an obstructed superficial femoral artery segment by anastomosing a graft (autologous great saphenous vein or PTFE) from the common femoral artery to the popliteal artery above or below the knee.

\medterm{Fogarty Catheter Embolectomy}-- Surgical technique for emergency extraction of arterial emboli or fresh thrombi; a balloon-tipped Fogarty catheter is passed distally beyond the occlusion, the balloon is inflated, and withdrawn, pulling the occlusive material out through arteriotomy.

\medterm{Giant Cell Arteritis (Temporal Arteritis)}-- Granulomatous vasculitis of large and medium-sized vessels affecting elderly adults ($>50\text{ years}$), targeting extracranial branches of the carotid artery (temporal artery). Symptoms: headache, jaw claudication, scalp tenderness, and amaurosis fugax. High risk of irreversible blindness due to anterior ischemic optic neuropathy; requires immediate high-dose corticosteroid therapy.

\medterm{Inferior Vena Cava (IVC) Filter}-- Metallic intravascular filter device deployed fluoroscopically into the infrarenal IVC to mechanically trap large emboli originating from lower extremity deep vein thrombosis, preventing pulmonary embolism. Indicated in patients with acute DVT/PE who have absolute contraindications to anticoagulation or recurrent PE despite adequate anticoagulation.

\medterm{Intermittent Claudication}-- Classic symptomatic manifestation of stable peripheral artery disease (PAD). Characterized by reproducible muscle aching, cramping, pain, or fatigue in the lower extremity (calf, thigh, or buttock) precipitated by physical exertion (walking) and promptly relieved by rest (usually within 2--5 minutes). Site of pain correlates with proximal arterial stenosis: calf claudication indicates superficial femoral or popliteal artery disease; thigh/buttock claudication indicates iliofemoral disease. Pathophysiology: exercise-induced imbalance between skeletal muscle oxygen demand and restricted blood supply. Management: risk factor modification (smoking cessation, statin, antiplatelet), structured exercise therapy, and cilostazol.

\medterm{Lymphedema}-- Accumulation of interstitial fluid and tissue swelling resulting from impaired lymphatic drainage, leading to chronic inflammation, adipose tissue deposition, and progressive subcutaneous fibrosis. Classified into Primary Lymphedema (congenital developmental malformations: Milroy disease, Meige disease) and Secondary Lymphedema (acquired damage to lymph nodes/vessels: surgical lymph node dissection for cancer, radiation therapy, tumor obstruction, or filariasis caused by Wuchereria bancrofti—the leading cause worldwide). Clinical features: non-pitting edema of the limb (Stemmer sign positive—inability to pinch a fold of skin at the base of the second toe), hyperkeratosis, and recurrent cellulitis. Management: Complete Decongestive Therapy (CDT: manual lymphatic drainage, multilayer compression bandaging, skin care, and exercise).

\medterm{May-Thurner Syndrome}-- Vascular compression anomaly where the left common iliac vein is chronically compressed against the lumbar spine by the overlying right common iliac artery, resulting in venous stasis, intraluminal fibrous bands, and a markedly increased risk of left iliofemoral DVT in young females.

\medterm{Open Aneurysm Repair}-- Traditional open surgical technique for abdominal or thoracic aortic aneurysm repair involving a laparotomy or thoracotomy, cross-clamping the aorta, opening the aneurysm sac, and sewing a synthetic prosthetic graft (Dacron or PTFE) in an inlay fashion.

\medterm{Paget-Schroetter Syndrome (Effort Thrombosis)}-- Acute primary deep vein thrombosis of the subclavian-axillary vein resulting from repetitive upper extremity mechanical compression between the clavicle and first rib (thoracic outlet obstruction), classically seen in young competitive athletes performing overhead arm motions (baseball pitchers, swimmers).

\medterm{Percutaneous Transluminal Angioplasty (PTA)}-- Endovascular interventional catheter procedure where a balloon catheter is advanced under fluoroscopy across a focal arterial stenosis and inflated to dilate the lumen by disrupting atherosclerotic plaque; often combined with metallic stent placement.

\medterm{Peripheral Artery Disease (PAD)}-- Systemic atherosclerotic occlusive disease of the arteries supplying the upper and lower extremities (excluding coronary and intracranial circulation). Affects over 200 million people worldwide. Major risk factors: cigarette smoking, diabetes mellitus, hypertension, hyperlipidemia, and advanced age. Clinical spectrum: asymptomatic (50\%), atypical leg pain, intermittent claudication, critical limb ischemia (rest pain, non-healing ulcers), and acute limb ischemia (the "6 Ps": Pain, Pallor, Pulselessness, Paresthesia, Paralysis, Poikilothermia). Diagnosis: Ankle-Brachial Index (ABI $\le 0.90$) and CTA/MRA. Treatment: medical therapy (aspirin/clopidogrel, high-intensity statin, ACE inhibitors), exercise therapy, and revascularization (angioplasty/stenting or surgical bypass).

\medterm{Popliteal Artery Aneurysm}-- The most common peripheral artery aneurysm ($70\text{--}80\%$), defined as popliteal artery diameter $>2\text{ cm}$. Strongly associated with abdominal aortic aneurysms ($50\%$ co-occurrence). Unlike AAAs which rupture, popliteal aneurysms present primarily with thromboembolism, causing acute limb ischemia or distal trash feet.

\medterm{Popliteal Artery Entrapment Syndrome}-- Developmental vascular anomaly where the popliteal artery is abnormally compressed by adjacent musculotendinous structures (medial head of gastrocnemius muscle) in the popliteal fossa, leading to exertional leg claudication in young, healthy active individuals.

\medterm{Post-Thrombotic Syndrome (PTS)}-- Chronic venous condition developing in up to $50\%$ of patients following deep vein thrombosis, caused by irreversible venous valve damage and residual venous obstruction. Symptoms include persistent lower extremity swelling, aching pain, heaviness, venous ectasia, hyperpigmentation, and stasis ulceration.

\medterm{Pulmonary Embolism (PE)}-- Occlusion of one or more pulmonary arteries by a thrombus (or less commonly fat, air, or amniotic fluid) that has dislodged from a distant site, overwhelmingly ($>90\%$) originating from a deep vein thrombosis (DVT) of the lower extremities. Pathophysiology: mechanical obstruction increases right ventricular afterload, risking acute right ventricular failure and cardiogenic shock. Symptoms: sudden dyspnea, pleuritic chest pain, cough, hemoptysis, tachypnea, tachycardia, and apprehension. Diagnosis: Well's score for probability, D-dimer, and CT Pulmonary Angiography (CTPA—gold standard). Management: anticoagulation (LMWH, DOACs), systemic thrombolysis (alteplase) for massive PE with hemodynamic instability, or catheter-directed embolectomy.

\medterm{Radiofrequency Ablation (Venous)}-- Minimally invasive endovascular technique delivering thermal energy generated by high-frequency alternating current via a catheter inserted into an incompetent saphenous vein, causing vein wall contraction, lumen obliteration, and venous reflux elimination.

\medterm{Raynaud Phenomenon}-- Exaggerated vasospastic response of digital arteries and arterioles to cold temperature or emotional stress, causing episodic digital ischemia. Classic triphasic color response: (1) White (pallor due to intense arterial vasospasm), (2) Blue (cyanosis due to deoxygenated blood accumulation), and (3) Red (rubor due to reactive hyperemia upon rewarming, accompanied by throbbing pain and paresthesias). Primary Raynaud Disease: benign, idiopathic, symmetrical, without underlying disease. Secondary Raynaud Syndrome: associated with underlying systemic autoimmune diseases (systemic sclerosis/scleroderma, lupus, rheumatoid arthritis), occupational vibration injury, or medications. Management: cold avoidance, keeping body warm, and dihydropyridine calcium channel blockers (nifedipine, amlodipine).

\medterm{Saphenous Vein Stripping}-- Surgical procedure for severe superficial venous insufficiency involving ligation of the saphenofemoral junction (flush ligation) and physical avulsion/removal of the great saphenous vein trunk from the leg using a surgical stripper wire.

\medterm{Subclavian Steal Syndrome}-- Hemodynamic phenomenon resulting from severe stenosis or occlusion of the proximal subclavian artery proximal to the vertebral artery origin. Exercise of the ipsilateral arm causes retrograde blood flow down the ipsilateral vertebral artery away from the brainstem, causing vertebrobasilar insufficiency (dizziness, ataxia, syncope) and arm claudication.

\medterm{Superior Vena Cava (SVC) Syndrome}-- Clinical emergency resulting from obstruction of blood flow through the superior vena cava, overwhelmingly ($80\text{--}90\%$) caused by malignant compression (small cell lung cancer, non-Hodgkin lymphoma). Features: facial edema, plethora, neck and upper chest vein distension, dyspnea, and headache.

\medterm{Thoracic Endovascular Aortic Repair (TEVAR)}-- Minimally invasive endovascular procedure deploying a flexible fabric-covered stent graft within the thoracic aorta via femoral artery access to treat thoracic aortic aneurysms, Type B aortic dissections, or traumatic aortic transections.

\medterm{Thoracic Outlet Syndrome (TOS)}-- Symptom complex resulting from compression of the neurovascular bundle (brachial plexus, subclavian artery, subclavian vein) at the thoracic outlet (interscalene triangle, costoclavicular space). Categorized as Neurogenic ($95\%$, pain/numbness in hand), Venous (Paget-Schroetter syndrome), or Arterial (subclavian aneurysm/ischemia).

\medterm{Thromboangiitis Obliterans}-- See Buerger Disease.

\medterm{Transcranial Doppler (TCD)}-- Non-invasive ultrasound modality measuring velocity and direction of blood flow within major basal intracranial arteries through thin skull windows (temporal, orbital); used to monitor cerebral vasospasm after subarachnoid hemorrhage and screen children with sickle cell disease for stroke risk.

\medterm{Ultrasound-Guided Foam Sclerotherapy}-- Minimally invasive venous treatment where a sclerosant liquid (polidocanol) mixed with air to form foam is injected into incompetent varicose veins under ultrasound guidance, displacing blood and producing chemical endovenous inflammation and sclerosis.

\medterm{Varicose Veins}-- Dilated, tortuous, elongated superficial veins $\ge 3\text{ mm}$ in diameter, most commonly affecting the great and small saphenous veins and their tributaries in the lower extremities. Caused by primary valvular incompetence, saphenofemoral junction incompetence, or weakness of the venous wall, leading to retrograde venous reflux and elevated hydrostatic pressure. Risk factors: female gender, pregnancy, prolonged standing, obesity, family history, and advanced age. Symptoms: aching, heaviness, leg fatigue, pruritus, cosmetic concern, and localized phlebitis. Treatment: compression hosiery, conservative measures, endovenous thermal ablation (radiofrequency or endovenous laser therapy—EVLT), ultrasound-guided foam sclerotherapy, or surgical ambulatory phlebectomy / vein stripping.

\medterm{Vascular Steal Syndrome}-- Complication of surgical arteriovenous fistula or graft created for hemodialysis, occurring when the low-resistance AV access diverts ('steals') arterial blood flow away from distal extremity tissue, causing distal hand ischemia, coldness, pain, and sensory loss.

\medterm{Vascular Surgery \& Angiology}-- Angiology is the medical subspecialty dedicated to the study, diagnosis, and non-surgical management of diseases of the circulatory system (arteries, veins, and lymphatic vessels). Vascular Surgery is the surgical specialty focused on open and endovascular interventions for structural, traumatic, inflammatory, and occlusive diseases of the vascular system (excluding intracranial and coronary vessels). Major clinical domains: abdominal/thoracic aortic aneurysms, peripheral artery disease, carotid artery stenosis, deep vein thrombosis, venous insufficiency, vascular access, and vascular trauma.

\medterm{Vein Graft Failure}-- Occlusion or stenosis of an autologous vein bypass graft caused early ($<30\text{ days}$) by technical error or thrombosis, mid-term ($1\text{ month to }2\text{ years}$) by intimal hyperplasia, or late ($>2\text{ years}$) by progressive atherosclerosis within graft wall.
